% ---------------------------------------------------------------------------
% CBMC — draft skeleton
%
% Compiles with plain pdflatex, no conference style files needed. Swap the
% documentclass for neurips_2026.sty / icml2026.sty when the venue is chosen.
%
% NOTE: the method section below was reconstructed from GOAL.md and the
% implemented code (cbmc/concepts/cem.py). Reconcile it against the formal
% framework document before submission — see \todo markers throughout.
% ---------------------------------------------------------------------------
\documentclass[11pt]{article}

\usepackage[margin=1in]{geometry}
\usepackage{amsmath,amssymb,amsthm}
\usepackage{booktabs}
\usepackage{graphicx}
\usepackage{multirow}
\usepackage{xcolor}
\usepackage{hyperref}
\usepackage{natbib}

\newcommand{\todo}[1]{\textcolor{red}{[TODO: #1]}}
\newcommand{\cvec}{\mathbf{c}}
\newcommand{\xvec}{\mathbf{x}}

\title{Continuous Concept Bottleneck Models:\\
       Interpolating Concept Embeddings over Real-Valued Concepts}
\author{\todo{authors} \\ Universit\`a della Svizzera italiana}
\date{}

\begin{document}
\maketitle

% ===========================================================================
\begin{abstract}
Concept Bottleneck Models (CBMs) make predictions interpretable by routing them
through a layer of human-meaningful concepts, but they assume concepts are
binary. Many concepts of practical interest are not: the diameter of a lesion,
the angle of a joint, the duration of a cardiac interval. Forcing such
quantities into binary form discards the information that makes them useful.
We present \textbf{CBMC}, a concept bottleneck that operates on continuous
concepts. Building on Concept Embedding Models (CEM), we represent each concept
by two learned anchor embeddings and place the sample at a point \emph{between}
them determined by a predicted scalar, so that a concept's magnitude -- not
merely its presence -- is carried forward to the task head. The same
formulation admits test-time intervention on real-valued concepts.
On an arithmetic-MNIST benchmark the continuous bottleneck halves task error
relative to a binary CBM ($31.2$ vs.\ $57.7$ MSE) and reduces concept error by
$8\times$, while also halving the error of an unconstrained backbone. The gap
is starkest under intervention: supplying ground-truth concepts \emph{reduces}
our error to $13.5$ but \emph{raises} the binary CBM's to $133.6$, because a binary
bottleneck cannot express the value being supplied. We additionally find that
discretising a continuous concept into $k$ bins is competitive --- better still
under intervention --- at a parameter cost growing linearly in $k$, and that
gating a single embedding multiplicatively is substantially worse than
interpolating between two anchors. On LIDC-IDRI lung CT, however, where
concepts are radiologist annotations rather than generative factors, the
advantage does not reproduce: over 30 seeds no concept representation is
distinguishable from any other on balanced accuracy, nor from an unconstrained
backbone, and a scalar bottleneck intervenes significantly better than ours. We
report this negative transfer, and what we think drives it, alongside the
synthetic results.
\end{abstract}

% ===========================================================================
\section{Introduction}
\label{sec:intro}

\todo{This section is a scaffold. Replace with the narrative from the method
PDF once available.}

Deep models are accurate and opaque. Concept Bottleneck Models
\citep{koh2020concept} address the opacity by forcing predictions through an
intermediate layer whose units correspond to concepts a human has named, so a
prediction can be read as a function of quantities a domain expert recognises,
and can be corrected by editing them.

The standard formulation treats concepts as binary. \emph{Is there spiculation?
Is the wing brown?} For many domains this is the wrong type. A radiologist does
not record that a nodule \emph{has} a diameter; they record that it is
\SI{12.7}{\milli\metre}. Binarising such a concept forces an arbitrary
threshold, destroys ordering, and makes intervention coarse: one can assert
``large'' but not ``\SI{12.7}{\milli\metre}''.

Concept Embedding Models \citep{zarlenga2022concept} improved CBM accuracy by
giving each concept a pair of learned embeddings and mixing them according to
the predicted probability that the concept is active. We observe that this
mixing coefficient need not be a probability. If it is instead a normalised
\emph{magnitude}, the same architecture represents continuous concepts, and the
mixing weight becomes a position along a learned axis between two anchors.

\paragraph{Contributions.}
\begin{enumerate}
  \item A concept bottleneck for continuous concepts, obtained by reinterpreting
        the CEM mixing coefficient as a clamped linear function of a predicted
        real-valued concept (Section~\ref{sec:method}).
  \item Intervention on real-valued concepts, and a protocol for measuring it;
        the binary baseline is \emph{harmed} by correct interventions while
        ours is helped.
  \item A controlled comparison of four concept representations --- binary,
        categorical, binned and continuous --- under one backbone, isolating
        representation from architecture (Section~\ref{sec:binning}).
  \item An ablation separating two-anchor interpolation from multiplicative
        gating, including the observation that output normalisation makes a
        multiplicative gate invariant to concept magnitude.
  \item A reusable LIDC-IDRI benchmark: 2{,}628 nodule volumes with eleven
        annotated concepts of three different types, plus measured annotator
        disagreement; and a negative transfer result on it, where the
        advantage measured on synthetic data does not reproduce
        (Section~\ref{sec:lidc}).
\end{enumerate}

% ===========================================================================
\section{Related Work}
\label{sec:related}

\paragraph{Concept bottlenecks.}
\citet{koh2020concept} introduced CBMs with binary concepts and the
independent/sequential/joint training regimes we follow.
\todo{cite and position: Post-hoc CBM \citep{yuksekgonul2023posthoc} and
Label-free CBM \citep{oikarinen2023labelfree}, both of which attach a
bottleneck to a frozen foundation model; neither addresses continuous concepts.}

\paragraph{Concept embeddings.}
CEM \citep{zarlenga2022concept} is our direct predecessor and the baseline we
compare against most closely. \todo{Make explicit that CEM is the $n=2$,
binary-probability special case of our formulation.}

\paragraph{Uncertainty over concepts.}
Probabilistic CBM \citep{kim2023probabilistic} models each concept as a
distribution rather than a point, and is the closest existing work to our
Case~4. \todo{We do not yet compare empirically; see
Section~\ref{sec:limitations}.}

\paragraph{Continuous and structured concepts.}
\todo{Survey what exists. Candidates: concept whitening, CBMs on regression
targets, disentanglement work on dSprites/Shapes3D where generative factors are
literally continuous. This paragraph is the one a reviewer will check hardest
for novelty -- do not leave it thin.}

% ===========================================================================
\section{Method}
\label{sec:method}

\todo{Reconcile the whole section against the formal framework PDF. What
follows matches the implementation in \texttt{cbmc/concepts/cem.py}.}

\subsection{Setup}
Let $\xvec$ be an input, $\cvec \in \mathbb{R}^k$ a vector of $k$ concept
annotations, and $y$ a task target. A backbone $f$ maps $\xvec$ to a latent
$h = f(\xvec)$. A bottleneck maps $h$ to concept predictions
$\hat{\cvec}$ and to an embedding $z$ consumed by a task head $g$, so that
$\hat{y} = g(z)$. Training minimises
\begin{equation}
  \mathcal{L} = \mathcal{L}_{\text{task}}(\hat{y}, y)
              + \lambda\,\mathcal{L}_{\text{concept}}(\hat{\cvec}, \cvec).
  \label{eq:loss}
\end{equation}

\subsection{A progression of concept types}
We organise the design space by what a concept \emph{is}.

\paragraph{Case 1: binary (CEM).}
Each concept $i$ has anchors $e_i^{+}, e_i^{-} \in \mathbb{R}^m$ and a predicted
probability $p_i = \sigma(g_i(h))$. The embedding is the convex combination
$z_i = p_i e_i^{+} + (1 - p_i) e_i^{-}$.

\paragraph{Case 2: categorical.}
An $n$-state concept gets $n$ anchors and a softmax over them:
\begin{equation}
  p_i = \mathrm{softmax}\big(g_i([\phi_{i,1} \Vert \dots \Vert \phi_{i,n}])\big),
  \qquad
  z_i = \sum_{s=1}^{n} p_{i,s}\, \phi_{i,s}(h).
  \label{eq:case2}
\end{equation}
Following CEM, the state distribution is predicted from the \emph{concatenated}
anchors rather than from $h$ directly. Case~1 is the $n=2$ special case, so a
single implementation covers both. Concepts may carry different state counts,
which matters for LIDC: \texttt{calcification} is a 6-way nominal code and
\texttt{internalStructure} a 4-way one, while the remaining ratings have five
levels.

\paragraph{Case 3: discretised ranges.}
A continuous range split into $k$ uniform bins, with bin identity treated as a
categorical concept under Eq.~\ref{eq:case2}. Binning imposes a floor on
achievable concept error --- the residual from reconstructing a value as its
bin centre --- which we report alongside measured error so the two can be
separated (Section~\ref{sec:binning}).

\paragraph{Case 3.5: continuous via two anchors.}
\textbf{This is the method evaluated in this paper.} Rather than discretise, we
keep only the two extreme anchors and let a concept-specific regressor place
the sample between them. For concept $i$ with predicted value $\hat{c}_i$,
\begin{equation}
  w_i = \mathrm{clamp}\!\left(\tfrac{1}{2}\hat{c}_i + \tfrac{1}{2},\, 0,\, 1\right),
  \qquad
  z_i = w_i\, e_i^{+}(h) + (1 - w_i)\, e_i^{-}(h).
  \label{eq:case35}
\end{equation}
Concept annotations are normalised to $[-1, 1]$, so $w_i$ spans $[0,1]$ and the
anchors are reached exactly at the extremes of the observed range. Anchors are
produced by per-concept networks $e_i^{\pm}(h)$ rather than being free
parameters, following CEM.

We deliberately apply no squashing nonlinearity to $\hat{c}_i$ before the clamp.
A $\tanh$ would compress the attainable range of $w_i$ to roughly
$[0.12, 0.88]$, weakening both predictions and interventions; we verify this
empirically in Section~\ref{sec:ablation}.

\paragraph{Case 4: distributional.}
As bins $\to \infty$ the concept becomes a density $p(z \mid \xvec)$.
\todo{Theory only. Either implement for a head-to-head against Probabilistic
CBM, or state explicitly as future work -- see Section~\ref{sec:limitations}.}

\subsection{Intervention}
At test time any subset of concepts may be replaced by ground truth. With a
binary mask $\mathbf{m}$,
\begin{equation}
  \tilde{c}_i = m_i c_i + (1 - m_i)\hat{c}_i,
\end{equation}
and $\tilde{c}$ is substituted for $\hat{c}$ in Eq.~\ref{eq:case35}. Because
$w_i$ is monotone in $\tilde{c}_i$ and saturates exactly at the anchors,
intervening with an extreme value yields the corresponding anchor embedding --
a vector the head has seen during training.

\subsection{Alternatives we rule out}
\label{sec:alternatives}
Two simpler designs replace the anchor pair with a single embedding
$\phi_i(h)$ scaled by the concept:
\begin{align}
  \textbf{CEM-Tanh:}   \quad & z_i = \tanh(\hat{c}_i)\,\phi_i(h), \\
  \textbf{CEM-Linear:} \quad & z_i = \hat{c}_i\,\phi_i(h).
\end{align}
Both are strictly cheaper -- one embedding network per concept instead of two.
We find both substantially worse under intervention
(Section~\ref{sec:ablation}).

A subtlety: with the output \textsc{LayerNorm} used throughout, CEM-Linear is
\emph{invariant} to any global rescaling of $\hat{\cvec}$, since
$\mathrm{LN}(\alpha z) = \mathrm{LN}(z)$. Only the relative pattern across
concepts reaches the head. We therefore also report \textbf{CEM-Linear-Raw},
identical but without the normalisation, which isolates this effect.
Empirically, scaling every concept by $10$ moves a normalised CEM-Linear's
output by $5\times 10^{-4}$ while the unnormalised variant scales exactly
tenfold, confirming the invariance is complete rather than approximate.

\subsection{Balancing the two losses}
\label{sec:lossbalance}
Eq.~\ref{eq:loss} trades a task term against a concept term, and the two are
not automatically on the same scale. On ArithmeticMNIST the target is an
arithmetic result spanning $[0, 81]$, so the task MSE is of order $30$, while
concepts are normalised to $[-1,1]$ and their MSE is of order $0.3$. At
$\lambda = 1$ the concept gradient is roughly $95\times$ smaller and the concept
head is effectively unsupervised: concept accuracy stays at chance and every
concept representation looks equally poor.

Raising $\lambda$ to $100$ changes the picture completely --- a 10-bin model
moves from $13\%$ to $94\%$ concept accuracy --- and reverses the ranking
between representations. We therefore report $\lambda$ with every result, and
recommend checking concept accuracy against chance before drawing any
conclusion about a bottleneck. \todo{$\lambda$ was not tuned; only $1$, $10$
and $100$ were tried.}

% ===========================================================================
\section{Experimental Setup}
\label{sec:setup}

\paragraph{Datasets.}
\emph{ArithmeticMNIST}: two MNIST digits and an operator rendered on one canvas;
concepts are the two digit values, the target is the arithmetic result.
\emph{Pendulum}: rendered images of a pendulum under a light source; concepts
are pendulum angle and light angle, the target is shadow geometry.
\emph{LIDC-IDRI} \citep{armato2011lidc}: 1{,}018 chest CT scans annotated by up
to four radiologists each. We extract a $64^3$\,mm box around every annotated
nodule, resampled to $1$\,mm isotropic voxels so that a fixed voxel count means
a fixed physical size regardless of scanner --- native slice spacing varies
from $0.6$ to $5$\,mm, an $8.3\times$ spread that would otherwise make size
concepts unlearnable from pixels. This yields \textbf{2{,}628 nodules from 875
patients} (Table~\ref{tab:lidc}). Concepts are six radiologist ratings on a
1--5 scale plus three geometric measurements (diameter, volume, surface area)
computed from the traced contours; the target is malignancy, binarised by the
standard convention. Splits are by patient: $71\%$ of nodules share a patient
with another and $56\%$ of those groups are label-homogeneous, so a
nodule-level split would leak the scanner rather than test the model.
\todo{Training results pending.}

\paragraph{Baselines.}
An unconstrained backbone with no bottleneck (\textbf{Baseline}); a binary
concept bottleneck (\textbf{CBM}); and \textbf{CEM}.
\todo{Add Probabilistic CBM, Post-hoc CBM, Label-free CBM -- none are currently
run, and a reviewer will expect at least Probabilistic CBM.}

\paragraph{Metrics.}
Task MSE; concept MSE against normalised annotations; \emph{intervention MSE},
the task error when all concepts are replaced by ground truth; and mean absolute
error. For generative models we report test reconstruction loss and KL.
All numbers are mean $\pm$ standard deviation over 5 seeds.

% ===========================================================================
\section{Results}
\label{sec:results}

The tables below carry the main argument. Every result produced so far,
including the ones not discussed here, is collected in
Appendix~\ref{app:tables} with its measurement conditions stated, since those
differ across experiments.

\subsection{Prediction}

\begin{table}[t]
\centering
\caption{ArithmeticMNIST, 5 seeds, $\lambda = 1$. Lower is better throughout.
CEM (our Case~3.5 formulation) improves on both the unconstrained baseline and
the binary CBM on every metric, and is the only concept model whose
interventions help rather than hurt. See Section~\ref{sec:lossbalance}: these
numbers use $\lambda = 1$, which the binning study suggests under-trains every
concept head.}
\label{tab:mnist}
\begin{tabular}{lcccc}
\toprule
\textbf{Model} & \textbf{Task MSE} & \textbf{Concept MSE} &
\textbf{Intervention MSE} & \textbf{MAE} \\
\midrule
Baseline (no bottleneck) & $61.18 \pm 6.07$ & --- & --- & $4.97 \pm 0.30$ \\
CBM \citep{koh2020concept}       & $57.74 \pm 2.42$ & $2.763 \pm 0.189$ & $133.64 \pm 0.81$ & $4.99 \pm 0.18$ \\
\textbf{CEM (ours)}              & $\mathbf{31.18 \pm 1.27}$ & $\mathbf{0.344 \pm 0.008}$ & $\mathbf{13.45 \pm 2.58}$ & $\mathbf{3.17 \pm 0.05}$ \\
\midrule
\multicolumn{5}{l}{\emph{Ablations (single embedding, multiplicative gate)}} \\
CEM-Tanh                         & $41.31 \pm 2.40$ & $0.402 \pm 0.006$ & $37.35 \pm 0.36$ & $3.88 \pm 0.09$ \\
CEM-Linear                       & $41.59 \pm 2.85$ & $0.406 \pm 0.036$ & $33.80 \pm 3.62$ & $3.79 \pm 0.24$ \\
CEM-Linear-Raw                   & \todo{not run} & \todo{} & \todo{} & \todo{} \\
\bottomrule
\end{tabular}
\end{table}

On ArithmeticMNIST (Table~\ref{tab:mnist}) the continuous bottleneck is not
merely competitive with the unconstrained baseline -- it is \emph{better},
halving its error ($31.18$ vs.\ $61.18$), while the binary CBM barely improves
on it ($57.74$). Concept supervision acts as a useful inductive bias when the
concepts are genuinely predictive of the target, which for a sum of two digits
they are by construction. The single-embedding ablations land between the two
($41.31$, $41.59$), so both components -- having anchor embeddings at all, and
having two of them to interpolate between -- contribute.

The intervention result is the sharper one. Replacing predicted concepts with
ground truth \emph{reduces} CEM's error to $13.45$ but \emph{increases} the
binary CBM's to $133.64$, more than double its own uncorrected error. A binary
bottleneck cannot express ``this digit is a 7''; the intervention pushes the
representation somewhere the head was never trained on.

\begin{table}[t]
\centering
\caption{Pendulum regression. \todo{These numbers are reported for
completeness but are not yet trustworthy -- see Section~\ref{sec:limitations}.}}
\label{tab:pendulum}
\begin{tabular}{lcccc}
\toprule
\textbf{Model} & \textbf{Task MSE} & \textbf{Concept MSE} &
\textbf{Intervention MSE} & \textbf{MAE} \\
\midrule
Baseline & $\mathbf{203.71 \pm 1.53}$ & --- & --- & $\mathbf{1.940 \pm 0.215}$ \\
CBM      & $210.22 \pm 1.70$ & $0.0199 \pm 0.0004$ & $\mathbf{2.110 \pm 0.248}$ & $2.623 \pm 0.049$ \\
CEM      & $206.51 \pm 2.54$ & $\mathbf{0.0193 \pm 0.0018}$ & $7.815 \pm 1.987$ & $2.218 \pm 0.204$ \\
\bottomrule
\end{tabular}
\end{table}

\subsection{Discretising a continuous concept}
\label{sec:binning}

Section~\ref{sec:method} offers two ways to handle a continuous concept: bin it
and treat the bin as categorical (Case~3), or interpolate between two anchors
(Case~3.5). Table~\ref{tab:binned} compares them directly on ArithmeticMNIST,
where the concepts are digit values in $[1,9]$, holding the backbone, embedding
size, learning rate and seed fixed so the representation is the only
difference.

\begin{table}[t]
\centering
\caption{Binned concepts (Case~3) versus two-anchor interpolation (Case~3.5) on
ArithmeticMNIST, 30 epochs, $\lambda = 100$. Binning wins on every error metric
and improves monotonically with resolution, but the parameter cost grows
linearly in $k$: each concept needs $k$ anchor networks rather than $2$.}
\label{tab:binned}
\begin{tabular}{lcccrc}
\toprule
\textbf{Model} & \textbf{Task MSE} & \textbf{Concept MSE} &
\textbf{Interv.\ MSE} & \textbf{Params} & \textbf{Quant.\ floor} \\
\midrule
CEM, two anchors (Case 3.5) & $29.00$ & $0.0723$ & $16.98$ & $159$K & --- \\
\midrule
Binned, $k=5$   & $26.34$ & $0.0925$ & $7.45$ & $247$K & $0.0167$ \\
Binned, $k=10$  & $13.15$ & $0.0357$ & $2.58$ & $392$K & $0.0044$ \\
Binned, $k=20$  & $\mathbf{10.09}$ & $\mathbf{0.0313}$ & $\mathbf{2.19}$ &
                  $683$K & $0.0014$ \\
\bottomrule
\end{tabular}
\end{table}

Two things stand out. First, discretisation does not cost accuracy here --- it
improves it, by nearly $3\times$ on task MSE at $k=20$ and by $7.8\times$ on
intervention error. A categorical intervention substitutes a one-hot, landing
exactly on a learned anchor the head has seen in training, whereas a continuous
intervention sets a weight that may fall anywhere between two anchors.

Second, that accuracy is bought with parameters: $683$K at $k=20$ against
$159$K for two anchors, growing linearly in $k$ while two anchors give
unbounded resolution at fixed cost. The quantisation floor column shows the
binned models remain $8$--$22\times$ above the error that binning alone
imposes, so resolution is not yet the binding constraint at these bin counts.

\todo{Single seed. Re-run with 5 seeds before claiming the ordering; the
supporting script already supports \texttt{-{}-runs}.}

\subsection{Generation}

\begin{table}[t]
\centering
\caption{Generative setting (convolutional VAE backbone). Test reconstruction
loss and KL. On ArithmeticMNIST the expected interpretability--fidelity
trade-off appears: the unconstrained VAE reconstructs best, the binary CBM
worst, and CEM sits between while achieving by far the lowest concept error.}
\label{tab:gen}
\begin{tabular}{lccc}
\toprule
\textbf{Dataset} & \textbf{Model} & \textbf{Recon.\ loss} & \textbf{Concept MSE} \\
\midrule
\multirow{3}{*}{ArithmeticMNIST}
 & Baseline VAE & $\mathbf{374.89 \pm 0.51}$ & --- \\
 & CBM          & $662.76 \pm 13.78$ & $1.058 \pm 0.478$ \\
 & CEM          & $430.59 \pm 1.49$ & $\mathbf{0.070 \pm 0.004}$ \\
\midrule
\multirow{3}{*}{Pendulum}
 & Baseline VAE & $330.16 \pm 0.00$ & --- \\
 & CBM          & $0.475 \pm 0.101$ & $0.154 \pm 0.016$ \\
 & CEM          & $165.23 \pm 164.94$ & $\mathbf{0.027 \pm 0.006}$ \\
\bottomrule
\end{tabular}
\end{table}

\todo{The pendulum rows of Table~\ref{tab:gen} must not be published as they
stand. The baseline VAE has posterior-collapsed (KL $= 0.0000 \pm 0.0000$,
identical across all seeds), and CEM's reconstruction standard deviation
($164.94$) is as large as its mean, indicating that some seeds collapsed and
others did not.}

\subsection{Ablation: what makes interpolation work}
\label{sec:ablation}

Replacing the anchor pair with a single multiplicatively-gated embedding costs
little in task error but a great deal under intervention: $13.45 \to 37.35$
(Tanh) and $13.45 \to 33.80$ (Linear), a factor of $2.5$--$2.8$. CEM-Linear-Raw, which removes the output
normalisation, is implemented and included in the LIDC comparison; on
ArithmeticMNIST it remains \todo{to be run}. If removing the normalisation
recovers intervention performance, the mechanism is the scale invariance of
Section~\ref{sec:alternatives}; if it does not, the two-anchor structure is
doing the work and the claim is stronger.

\subsection{Scaling to clinical data: LIDC-IDRI}
\label{sec:lidc}

The benchmarks above are synthetic: concepts are known exactly because they were
used to generate the data. LIDC-IDRI tests whether the formulation survives
genuine clinical annotation, where concepts are noisy human judgements and the
target is a diagnosis.

\begin{table}[t]
\centering
\caption{LIDC-IDRI after preprocessing. Each nodule is a $64^3$\,mm box at
$1$\,mm isotropic resolution with eleven annotated concepts. Requiring three or
more radiologists halves the set but balances the classes and removes nodules
only one annotator even identified.}
\label{tab:lidc}
\begin{tabular}{lrr}
\toprule
\textbf{Filter} & \textbf{Nodules} & \textbf{Benign / malignant} \\
\midrule
all extracted                      & 2{,}628 & --- \\
drop ambiguous (malignancy $=3$)   & 1{,}991 & $68.2\% / 31.8\%$ \\
\quad + at least 3 radiologists    & 1{,}180 & $55.7\% / 44.3\%$ \\
\bottomrule
\end{tabular}
\end{table}

\paragraph{Both concept types in one dataset.}
LIDC exercises the full progression. Diameter, volume and surface area are
genuine continuous scalars in physical units (Case~3.5). The six perceptual
ratings are ordinal 1--5 levels (Case~2). Two further annotations ---
\texttt{calcification} and \texttt{internalStructure} --- are nominal codes
rather than magnitudes, and we exclude them from continuous treatment; they are
the natural categorical case.

\paragraph{Annotator disagreement as a target for Case 4.}
Because up to four radiologists rated each nodule independently, every concept
carries a measured spread. Across nodules with at least two annotators the mean
standard deviation is $0.44$--$0.61$ rating points on a 1--5 scale, and
$10$--$14\%$ of nodules show strong disagreement ($\sigma > 1$). This is
ground-truth human uncertainty, not a modelling assumption, and gives Case~4 a
target to be evaluated against rather than merely asserted --- a stronger
evaluation than Probabilistic CBM \citep{kim2023probabilistic} currently
admits. \todo{Case 4 is still unimplemented; this is the argument for building
it.}

\paragraph{Seven representations, one encoder.}
Table~\ref{tab:lidcresults} compares seven concept representations under an
identical 3D CNN encoder, split, schedule and seeds. The result does not
support the ordering found on the synthetic benchmarks, and we report it as
measured.

\begin{table}[t]
\centering
\caption{LIDC-IDRI malignancy classification, \textbf{30 seeds}, 100 epochs,
$\lambda = 1$. Classes are $56/44$, so balanced accuracy is the figure to read.
``Interv.'' supplies the true concepts at test time. No pairwise difference in
balanced accuracy is statistically significant (Welch, $n=30$, all
$|t| < 2$), including against the unconstrained backbone. The intervention
column does separate: CBM is significantly better than ours ($t = 5.5$), as is
the categorical variant ($t = 2.4$).}
\label{tab:lidcresults}
\begin{tabular}{lcccr}
\toprule
\textbf{Representation} & \textbf{Balanced acc.} & \textbf{Interv.\ acc.} &
\textbf{\% of ceiling} & \textbf{Params} \\
\midrule
No bottleneck            & $0.778 \pm 0.063$ & --- & --- & $292$K \\
\midrule
CBM (scalar)             & $0.778 \pm 0.042$ & $\mathbf{0.846 \pm 0.015}$ & $98\%$ & $294$K \\
CEM binary (Case 1)      & $\mathbf{0.787 \pm 0.037}$ & $0.799 \pm 0.041$ & $93\%$ & $591$K \\
Categorical (Case 2/3)   & $0.782 \pm 0.038$ & $0.821 \pm 0.035$ & $95\%$ & $985$K \\
CEM-Linear (normalised)  & $0.778 \pm 0.036$ & $0.815 \pm 0.035$ & $94\%$ & $461$K \\
CEM-Linear (raw)         & $0.737 \pm 0.054$ & $0.758 \pm 0.050$ & $88\%$ & $461$K \\
CEM (ours, Case 3.5)     & $0.765 \pm 0.057$ & $0.795 \pm 0.049$ & $92\%$ & $591$K \\
\bottomrule
\end{tabular}
\end{table}

\paragraph{On this dataset the bottleneck is free, and our formulation is not
the best one.}
Every representation except CEM-Linear with raw values lands between $0.765$
and $0.787$ balanced accuracy, and so does the unconstrained backbone at
$0.778$. With 30 seeds none of those differences reach significance. Two
readings follow, and both matter.

First, the bottleneck costs nothing here: constraining predictions through nine
interpretable concepts is as accurate as not constraining them at all. That is
a positive result for concept bottlenecks in general.

Second, it is not a result for \emph{our} formulation specifically. Two-anchor
interpolation ranks sixth of seven on balanced accuracy, below binary
concepts ($0.787$), categorical ($0.782$) and a scalar CBM ($0.778$), though
again not significantly so. The advantage we measure on ArithmeticMNIST
(Table~\ref{tab:mnist}) does not reproduce when the concepts are radiologist
annotations rather than the generative factors of the data.

\paragraph{Intervention: the clearest signal, and it is against us.}
Unlike balanced accuracy, intervention accuracy does separate the methods, and
the scalar CBM wins decisively: $0.846$ against our $0.795$, a difference of
$5.5$ standard errors. The categorical variant also beats us significantly
($0.821$, $t = 2.4$). The mechanism is visible once the concepts-only ceiling
is in view (Table~\ref{tab:app-ceiling}): intervening on a scalar bottleneck
replaces its $k$ numbers outright, so the intervened model \emph{is} the
concepts-only classifier, and reaches $98\%$ of its $0.863$. An embedding
bottleneck must route the supplied value through anchor embeddings, and on this
data that indirection costs more than it returns.

This directly contradicts the ArithmeticMNIST finding, where interventions
harmed the binary CBM and helped ours. We do not have an account that explains
both. \todo{The partial-intervention curve --- accuracy against the fraction of
concepts corrected --- is the obvious next measurement, and would show whether
the embedding bottleneck has any regime where it leads.}

\paragraph{What does replicate: normalisation.}
CEM-Linear with raw concept values is last on both metrics ($0.737$, $0.758$)
and is the only variant significantly worse than ours under intervention
($t = 2.9$). With concepts spanning $[0, 10^4]$ in raw units, a multiplicative
gate is dominated by whichever concept has the largest magnitude, exactly as
Section~\ref{sec:alternatives} predicts.

\paragraph{Every model overfits.}
Training accuracy reaches $0.96$--$1.00$ against test balanced accuracy of
$0.74$--$0.79$ (Table~\ref{tab:app-overfit}), a gap of roughly $20$ points on
$944$ training volumes. No validation split or early stopping was used, so each
number is test performance at an arbitrary stopping point rather than at the
model's best. The single exception is CEM-Linear with raw values, which reaches
only $0.771$ training accuracy --- it fails to fit at all, consistent with the
magnitude-dominated gate. Augmentation was available but not enabled for this
run. We regard the LIDC comparison as under-regularised and therefore
preliminary. \todo{Re-run with augmentation, weight decay and early stopping on
a held-out validation split before drawing conclusions from the ordering.}

% ===========================================================================
\section{Limitations}
\label{sec:limitations}

\todo{Be explicit about all of the following; a reviewer will find them.}
\begin{itemize}
  \item \textbf{Case 4 is theory only.} The comparison against Probabilistic
        CBM \citep{kim2023probabilistic} that a reviewer will ask for first has
        not been run.
  \item \textbf{The pendulum benchmark is degenerate.} Its concept vector and
        task target are the same four physical quantities, so intervening with
        ground-truth concepts hands the model the answer. This explains why the
        binary CBM appears to beat CEM there ($0.473$ vs.\ $8.014$): the task
        rewards copying the input to the output, which a scalar bottleneck does
        more faithfully than an embedding blend. Either re-scope the task
        (concepts = angles, target = shadow geometry) or drop the benchmark.
  \item \textbf{Pendulum error is reported in raw physical units} and is
        dominated by one concept whose variance is 96\% of the total, making the
        aggregate number hard to compare against published results.
  \item \textbf{Posterior collapse} in the generative pendulum baseline
        (Table~\ref{tab:gen}).
  \item \textbf{Our formulation does not win on LIDC.} It ranks sixth of
        seven on balanced accuracy and is significantly worse than a scalar
        CBM under intervention ($t = 5.5$). The method's advantage is
        currently demonstrated only on data whose concepts are the generative
        factors.
  \item \textbf{Nothing is significant on LIDC balanced accuracy.} Over 30
        seeds no representation separates from any other, nor from the
        unconstrained backbone. The dataset may simply lack the power to rank
        these methods at $1{,}180$ samples.
  \item \textbf{The LIDC run is under-regularised.} Every model reaches
        $0.96$--$1.00$ training accuracy against $0.74$--$0.79$ test
        (Table~\ref{tab:app-overfit}). No augmentation, weight decay, early
        stopping or validation split was used, so the reported numbers are
        test performance at an arbitrary epoch.
  \item \textbf{One dataset, one split protocol}, with no external validation
        cohort.
  \item \textbf{Single-seed results.} Table~\ref{tab:binned} is one seed.
  \item \textbf{$\lambda$ untuned.} Section~\ref{sec:lossbalance} shows the
        conclusion about which representation is best can invert with the
        concept-loss weight. Only three values were tried, and the main
        ArithmeticMNIST results in Table~\ref{tab:mnist} use $\lambda = 1$,
        which the binning study suggests is too low --- those numbers may
        understate every concept model.
\end{itemize}

% ===========================================================================
\section{Conclusion}
\todo{Write last.}


% ===========================================================================
\appendix
\section{All experimental results}
\label{app:tables}

Every number produced so far, in one place. Each caption states the conditions
it was measured under, because they differ: the ArithmeticMNIST and pendulum
runs use $\lambda = 1$ and 5 seeds, the binning study uses $\lambda = 100$ and
a single seed, and the LIDC rows are dataset statistics and a concepts-only
ceiling rather than trained bottlenecks.

\subsection{ArithmeticMNIST --- prediction}

\begin{table}[h]
\centering
\caption{ArithmeticMNIST, 5 seeds, 20 epochs, $\lambda = 1$, intervention
probability $0.25$. Lower is better. ``Interv.'' replaces both predicted
concepts with ground truth.}
\label{tab:app-mnist}
\begin{tabular}{lcccc}
\toprule
\textbf{Model} & \textbf{Task MSE} & \textbf{Concept MSE} &
\textbf{Interv.\ MSE} & \textbf{MAE} \\
\midrule
Baseline (no bottleneck) & $61.18 \pm 6.07$ & --- & --- & $4.97 \pm 0.30$ \\
CBM                      & $57.74 \pm 2.42$ & $2.763 \pm 0.189$ & $133.64 \pm 0.81$ & $4.99 \pm 0.18$ \\
CEM (ours)               & $\mathbf{31.18 \pm 1.27}$ & $\mathbf{0.344 \pm 0.008}$ & $\mathbf{13.45 \pm 2.58}$ & $\mathbf{3.17 \pm 0.05}$ \\
CEM-Tanh                 & $41.31 \pm 2.40$ & $0.402 \pm 0.006$ & $37.35 \pm 0.36$ & $3.88 \pm 0.09$ \\
CEM-Linear               & $41.59 \pm 2.85$ & $0.406 \pm 0.036$ & $33.80 \pm 3.62$ & $3.79 \pm 0.24$ \\
\bottomrule
\end{tabular}
\end{table}

\begin{table}[h]
\centering
\caption{Two restricted variants of the same ArithmeticMNIST task, 5 seeds,
20 epochs, $\lambda = 1$. \texttt{all\_digits} uses digits 1--9;
\texttt{digit9} restricts to a single digit, which makes the task nearly
trivial and inflates every ratio. The ordering between models is unchanged.}
\label{tab:app-mnist-variants}
\begin{tabular}{llcccc}
\toprule
\textbf{Variant} & \textbf{Model} & \textbf{Task MSE} & \textbf{Concept MSE} &
\textbf{Interv.\ MSE} & \textbf{MAE} \\
\midrule
\multirow{5}{*}{\texttt{all\_digits}}
 & Baseline   & $71.63 \pm 5.93$ & --- & --- & --- \\
 & CBM        & $68.12 \pm 7.70$ & $3.871 \pm 3.341$ & $137.25 \pm 4.63$ & $5.45 \pm 0.32$ \\
 & CEM        & $\mathbf{46.72 \pm 4.03}$ & $\mathbf{0.418 \pm 0.027}$ & $\mathbf{23.70 \pm 4.13}$ & $\mathbf{4.10 \pm 0.20}$ \\
 & CEM-Tanh   & $63.80 \pm 9.85$ & $0.484 \pm 0.052$ & $57.10 \pm 18.48$ & $4.99 \pm 0.63$ \\
 & CEM-Linear & $56.32 \pm 9.96$ & $0.467 \pm 0.034$ & $49.51 \pm 11.18$ & $4.58 \pm 0.44$ \\
\midrule
\multirow{5}{*}{\texttt{digit9}}
 & Baseline   & $0.925 \pm 0.440$ & --- & --- & --- \\
 & CBM        & $2.302 \pm 2.636$ & $4.153 \pm 5.229$ & $992.40 \pm 0.66$ & $0.943 \pm 0.831$ \\
 & CEM        & $\mathbf{0.018 \pm 0.006}$ & $\mathbf{0.013 \pm 0.014}$ & $\mathbf{0.019 \pm 0.008}$ & $0.130 \pm 0.024$ \\
 & CEM-Tanh   & $0.085 \pm 0.077$ & $0.768 \pm 0.680$ & $14.42 \pm 14.77$ & $0.231 \pm 0.083$ \\
 & CEM-Linear & $0.045 \pm 0.081$ & $0.744 \pm 0.505$ & $8.94 \pm 9.57$ & $\mathbf{0.127 \pm 0.139}$ \\
\bottomrule
\end{tabular}
\end{table}

\subsection{Pendulum --- regression}

\begin{table}[h]
\centering
\caption{Pendulum, 5 seeds, 10 epochs, $\lambda = 1$. Errors are in raw
physical units, dominated by one concept whose variance is $96\%$ of the total.
\todo{This benchmark is degenerate --- concepts and target are the same four
quantities --- so intervention measures copying, not correction. See
Section~\ref{sec:limitations}.}}
\label{tab:app-pendulum}
\begin{tabular}{lcccc}
\toprule
\textbf{Model} & \textbf{Task MSE} & \textbf{Concept MSE} &
\textbf{Interv.\ MSE} & \textbf{MAE} \\
\midrule
Baseline & $\mathbf{203.71 \pm 1.53}$ & --- & --- & $\mathbf{1.940 \pm 0.215}$ \\
CBM      & $210.22 \pm 1.70$ & $0.0199 \pm 0.0004$ & $\mathbf{2.110 \pm 0.248}$ & $2.623 \pm 0.049$ \\
CEM      & $206.51 \pm 2.54$ & $\mathbf{0.0193 \pm 0.0018}$ & $7.815 \pm 1.987$ & $2.218 \pm 0.204$ \\
\bottomrule
\end{tabular}
\end{table}

\subsection{Generation (convolutional VAE backbone)}

\begin{table}[h]
\centering
\caption{Generative setting, 5 seeds, 20 epochs. Test reconstruction loss and
KL. \todo{The pendulum block is not reportable: the baseline VAE has
posterior-collapsed (KL identically $0$ across all seeds) and CEM's
reconstruction standard deviation equals its mean.}}
\label{tab:app-gen}
\begin{tabular}{llccc}
\toprule
\textbf{Dataset} & \textbf{Model} & \textbf{Recon.\ loss} & \textbf{KL} &
\textbf{Concept MSE} \\
\midrule
\multirow{3}{*}{ArithmeticMNIST}
 & Baseline VAE & $\mathbf{374.89 \pm 0.51}$ & $50.18 \pm 0.14$ & --- \\
 & CBM          & $662.76 \pm 13.78$ & $9.77 \pm 0.24$ & $1.058 \pm 0.478$ \\
 & CEM          & $430.59 \pm 1.49$ & $47.23 \pm 0.45$ & $\mathbf{0.070 \pm 0.004}$ \\
\midrule
\multirow{3}{*}{Pendulum}
 & Baseline VAE & $330.16 \pm 0.00$ & $0.00 \pm 0.00$ & --- \\
 & CBM          & $0.475 \pm 0.101$ & $77.27 \pm 9.09$ & $0.154 \pm 0.016$ \\
 & CEM          & $165.23 \pm 164.94$ & $46.46 \pm 36.04$ & $\mathbf{0.027 \pm 0.006}$ \\
\bottomrule
\end{tabular}
\end{table}

\subsection{Binned concepts --- effect of the concept-loss weight}

\begin{table}[h]
\centering
\caption{The same binning sweep at two values of $\lambda$, 30 epochs, 1 seed.
At $\lambda = 1$ the concept head is roughly $95\times$ under-weighted and no
model learns the concepts: every concept accuracy is at chance ($1/k$), and all
three binned models score worse than a predictor that ignores the image
entirely ($0.4167$). At $\lambda = 100$ the ranking inverts and binning wins
throughout. This table is the reason $\lambda$ is reported everywhere.}
\label{tab:app-lambda}
\begin{tabular}{lcccccc}
\toprule
& \multicolumn{3}{c}{$\lambda = 1$} & \multicolumn{3}{c}{$\lambda = 100$} \\
\cmidrule(lr){2-4}\cmidrule(lr){5-7}
\textbf{Model} & \textbf{Task} & \textbf{Concept} & \textbf{Conc.\ acc} &
                 \textbf{Task} & \textbf{Concept} & \textbf{Conc.\ acc} \\
\midrule
CEM (Case 3.5) & $32.44$ & $0.3424$ & $0.140$ & $29.00$ & $0.0723$ & $0.426$ \\
Binned $k=5$   & $33.97$ & $0.9346$ & $0.231$ & $26.34$ & $0.0925$ & $0.897$ \\
Binned $k=10$  & $36.63$ & $0.6387$ & $0.131$ & $13.15$ & $0.0357$ & $0.944$ \\
Binned $k=20$  & $35.07$ & $0.7098$ & $0.123$ & $\mathbf{10.09}$ & $\mathbf{0.0313}$ & $\mathbf{0.948}$ \\
\bottomrule
\end{tabular}
\end{table}

\begin{table}[h]
\centering
\caption{Binning at $\lambda = 100$ in full, 30 epochs, 1 seed. The
quantisation floor is the concept MSE a perfect binned model could not beat,
given that a value is reconstructed as its bin centre. All binned models sit
$8$--$22\times$ above their own floor, so resolution is not the binding
constraint at these bin counts.}
\label{tab:app-binned}
\begin{tabular}{lccccrc}
\toprule
\textbf{Model} & \textbf{Task MSE} & \textbf{Concept MSE} &
\textbf{Interv.\ MSE} & \textbf{Conc.\ acc} & \textbf{Params} &
\textbf{Quant.\ floor} \\
\midrule
CEM, two anchors & $29.00$ & $0.0723$ & $16.98$ & $0.426$ & $159$K & --- \\
Binned $k=5$     & $26.34$ & $0.0925$ & $7.45$  & $0.897$ & $247$K & $0.0167$ \\
Binned $k=10$    & $13.15$ & $0.0357$ & $2.58$  & $0.944$ & $392$K & $0.0044$ \\
Binned $k=20$    & $\mathbf{10.09}$ & $\mathbf{0.0313}$ & $\mathbf{2.19}$ &
                   $\mathbf{0.948}$ & $683$K & $0.0014$ \\
\bottomrule
\end{tabular}
\end{table}

\subsection{LIDC-IDRI --- dataset}

\begin{table}[h]
\centering
\caption{LIDC-IDRI after preprocessing: a $64^3$\,mm box per nodule at $1$\,mm
isotropic resolution, cropped from the full scan and the scan then discarded
(1.35\,GB kept rather than 128\,GB). Malignancy is averaged over annotators and
binarised at $3$; nodules landing exactly on $3$ are ambiguous and dropped.}
\label{tab:app-lidc-data}
\begin{tabular}{lrrr}
\toprule
\textbf{Filter} & \textbf{Nodules} & \textbf{Patients} & \textbf{Benign / malignant} \\
\midrule
all extracted                    & 2{,}628 & 875 & --- \\
drop ambiguous                   & 1{,}991 & --- & $68.2\% / 31.8\%$ \\
\quad + at least 3 radiologists  & 1{,}180 & 638 & $55.7\% / 44.3\%$ \\
\bottomrule
\end{tabular}
\end{table}

\begin{table}[h]
\centering
\caption{Annotator disagreement on LIDC, over nodules rated by more than one
radiologist: mean standard deviation of the rating, the fraction where they
agreed exactly, and the fraction differing by more than one point. This is
measured human uncertainty, available as an evaluation target for Case~4.}
\label{tab:app-disagreement}
\begin{tabular}{lccc}
\toprule
\textbf{Concept} & \textbf{Mean $\sigma$} & \textbf{$\sigma = 0$} & \textbf{$\sigma > 1$} \\
\midrule
subtlety    & $0.605$ & $19.5\%$ & $12.1\%$ \\
sphericity  & $0.573$ & $15.1\%$ & $6.3\%$ \\
margin      & $0.586$ & $20.6\%$ & $13.7\%$ \\
lobulation  & $0.511$ & $32.0\%$ & $12.2\%$ \\
spiculation & $0.437$ & $40.9\%$ & $10.4\%$ \\
texture     & $0.346$ & $53.7\%$ & $7.5\%$ \\
\midrule
malignancy (target) & $0.567$ & $21.8\%$ & $10.8\%$ \\
\bottomrule
\end{tabular}
\end{table}

\subsection{LIDC-IDRI --- concepts-only ceiling}

\begin{table}[h]
\centering
\caption{Classifiers given the annotated concepts and \emph{no image},
predicting malignancy. 5 seeds, patient-level split reseeded each run, 1{,}180
nodules. This bounds any concept bottleneck on this dataset from above: a
bottleneck routes every prediction through these concepts, so with a perfect
extractor it could match this and no more. It is also the value intervention
accuracy should approach, since intervening supplies exactly these values.
Tree models are scale-invariant and identical under both scalings, as expected.}
\label{tab:app-ceiling}
\begin{tabular}{llccc}
\toprule
\textbf{Model} & \textbf{Concepts} & \textbf{Accuracy} &
\textbf{Balanced acc.} & \textbf{AUC} \\
\midrule
\multirow{2}{*}{XGBoost}
 & raw        & $0.8528 \pm 0.007$ & $0.8500 \pm 0.007$ & $0.9241 \pm 0.008$ \\
 & normalised & $0.8528 \pm 0.007$ & $0.8500 \pm 0.007$ & $0.9241 \pm 0.008$ \\
\multirow{2}{*}{Random forest}
 & raw        & $\mathbf{0.8656 \pm 0.014}$ & $\mathbf{0.8628 \pm 0.015}$ & $0.9219 \pm 0.007$ \\
 & normalised & $\mathbf{0.8656 \pm 0.014}$ & $\mathbf{0.8628 \pm 0.015}$ & $0.9219 \pm 0.008$ \\
\multirow{2}{*}{Logistic regression}
 & raw        & $0.8502 \pm 0.006$ & $0.8485 \pm 0.006$ & $\mathbf{0.9247 \pm 0.009}$ \\
 & normalised & $0.8427 \pm 0.014$ & $0.8399 \pm 0.014$ & $0.9232 \pm 0.009$ \\
\multirow{2}{*}{MLP (64, 32)}
 & raw        & $0.8188 \pm 0.021$ & $0.8122 \pm 0.021$ & $0.8805 \pm 0.037$ \\
 & normalised & $0.8259 \pm 0.011$ & $0.8192 \pm 0.012$ & $0.9066 \pm 0.008$ \\
\bottomrule
\end{tabular}
\end{table}

\begin{table}[h]
\centering
\caption{XGBoost feature importance on the concepts-only task. The three
geometric concepts --- the genuinely continuous ones --- carry $54\%$ of the
signal between them, which is the case for treating concepts as real-valued
rather than categorical on this dataset.}
\label{tab:app-importance}
\begin{tabular}{lc}
\toprule
\textbf{Concept} & \textbf{Importance} \\
\midrule
diameter      & $0.277$ \\
volume        & $0.150$ \\
surface area  & $0.109$ \\
\midrule
texture       & $0.086$ \\
spiculation   & $0.085$ \\
lobulation    & $0.084$ \\
subtlety      & $0.073$ \\
margin        & $0.073$ \\
sphericity    & $0.062$ \\
\bottomrule
\end{tabular}
\end{table}

\subsection{LIDC-IDRI --- full comparison}

\begin{table}[h]
\centering
\caption{Every metric from the seven-variant LIDC run: 30 seeds, 100 epochs,
$\lambda = 1$, intervention probability $0.25$, encoder $[16,32,64,128]$,
embedding dimension $16$, $1{,}180$ nodules at $\geq 3$ annotators, no
augmentation. Concept MAE is relative to each concept's range, averaged over
the nine concepts; concept accuracy applies only to representations predicting
a discrete state.}
\label{tab:app-lidc-full}
\begin{tabular}{lccccc}
\toprule
\textbf{Representation} & \textbf{Accuracy} & \textbf{Balanced acc.} &
\textbf{Interv.\ acc.} & \textbf{Conc.\ acc.} & \textbf{Conc.\ MAE (rel)} \\
\midrule
No bottleneck    & $0.779 \pm 0.084$ & $0.778 \pm 0.063$ & --- & --- & --- \\
CBM (scalar)     & $0.784 \pm 0.047$ & $0.778 \pm 0.042$ & $\mathbf{0.846 \pm 0.015}$ & --- & $0.140$ \\
CEM binary       & $\mathbf{0.792 \pm 0.041}$ & $\mathbf{0.787 \pm 0.037}$ & $0.799 \pm 0.041$ & $0.744$ & $0.137$ \\
Categorical      & $0.788 \pm 0.043$ & $0.782 \pm 0.038$ & $0.821 \pm 0.035$ & $0.678$ & $0.140$ \\
CEM-Linear (norm)& $0.788 \pm 0.038$ & $0.778 \pm 0.036$ & $0.815 \pm 0.035$ & --- & $0.129$ \\
CEM-Linear (raw) & $0.744 \pm 0.067$ & $0.737 \pm 0.054$ & $0.758 \pm 0.050$ & --- & $\mathbf{0.120}$ \\
CEM (ours)       & $0.774 \pm 0.048$ & $0.765 \pm 0.057$ & $0.795 \pm 0.049$ & --- & $0.128$ \\
\bottomrule
\end{tabular}
\end{table}

Note again that CEM-Linear (raw) attains the \emph{lowest} concept error while
giving the \emph{worst} task and intervention accuracy. Predicting the concepts
and using them are separable failures: with raw magnitudes the multiplicative
gate is dominated by one concept, so the head largely ignores the bottleneck
even though the concept head itself trains well. Concept accuracy alone does
not make a bottleneck useful.

\begin{table}[h]
\centering
\caption{Overfitting on LIDC. Training accuracy is the mean over 30 seeds at
epoch 100; test balanced accuracy is from Table~\ref{tab:app-lidc-full}. With
$944$ training volumes and no augmentation, weight decay or early stopping,
every model that fits the data memorises it. CEM-Linear (raw) is the exception
only because it fails to fit.}
\label{tab:app-overfit}
\begin{tabular}{lccc}
\toprule
\textbf{Representation} & \textbf{Train acc.} & \textbf{Test balanced acc.} &
\textbf{Gap} \\
\midrule
No bottleneck     & $0.997$ & $0.778$ & $0.219$ \\
CEM binary        & $0.987$ & $0.787$ & $0.200$ \\
Categorical       & $0.987$ & $0.782$ & $0.205$ \\
CEM-Linear (norm) & $0.985$ & $0.778$ & $0.207$ \\
CEM (ours)        & $0.982$ & $0.765$ & $0.217$ \\
CBM (scalar)      & $0.957$ & $0.778$ & $0.179$ \\
CEM-Linear (raw)  & $0.771$ & $0.737$ & $0.034$ \\
\bottomrule
\end{tabular}
\end{table}

\subsection{Model sizes}\subsection{Model sizes}

\begin{table}[h]
\centering
\caption{Parameter counts for the LIDC comparison, all sharing the same
$[16,32,64,128]$ 3D CNN encoder with $9$ concepts and embedding dimension $16$.
The bottleneck is what differs. Categorical is the most expensive because each
concept needs one anchor network per state.}
\label{tab:app-sizes}
\begin{tabular}{lrr}
\toprule
\textbf{Variant} & \textbf{Parameters} & \textbf{Size (MB)} \\
\midrule
No bottleneck            & $291{,}714$ & $1.17$ \\
CBM (scalar)             & $294{,}059$ & $1.18$ \\
CEM-Linear               & $460{,}955$ & $1.84$ \\
CEM / CEM binary         & $591{,}275$ & $2.37$ \\
Categorical (5 states)   & $984{,}575$ & $3.94$ \\
\bottomrule
\end{tabular}
\end{table}

\todo{Still to fill: the seven-variant LIDC comparison (pipeline verified, run
not executed); CEM-Linear-Raw on ArithmeticMNIST; multi-seed repeats of
Tables~\ref{tab:app-lambda} and~\ref{tab:app-binned}; any Case~4 result.}

\bibliographystyle{plainnat}
\bibliography{refs}

\end{document}
